\section{Gap Analysis \& Related Work}
To justify the necessity of the Scriptorium Framework, we identify three critical gaps in contemporary hermeneutical software and pedagogical literature. Table \ref{tab:gaps} summarizes how the proposed framework addresses these deficiencies.

\begin{table*}[t]
\caption{Gap Analysis: Existing Methods vs. Scriptorium Framework}
\label{tab:gaps}
\centering
\begin{tabular}{@{}p{0.3\textwidth}p{0.3\textwidth}p{0.3\textwidth}@{}}
\toprule
\textbf{Identified Gap in Literature} & \textbf{Traditional AI Limitations} & \textbf{Scriptorium Improvement} \\ \midrule
\textbf{The Methodological Trap:} Mechanical checklists replace relational intimacy \cite{traina1952}. & Default to rigid ``fill-in-the-blank'' structures that mirror the OIA steps without depth. & \textbf{Method as Servant:} AI uses technical depth (Semantic Domains) to fuel natural conversation, not enforce a checklist. \\ \midrule
\textbf{Subjective Projection:} Personal trauma or culture projected onto the ``Father'' metaphor. & Mirroring of user bias or lack of historical correction \cite{bender2021parrots}. & \textbf{The Socratic Mirror:} Layer II (Historical Shock) uses AI to distinguish the biblical \textit{patria} from cultural projection \cite{malina2001}. \\ \midrule
\textbf{Rigor-Intimacy Disconnect:} Academic works (e.g., \cite{osborne2006}) are inaccessible; devotional guides are ``thin.'' & Provide either dry lexical data or superficial encouragement. & \textbf{The Computational Bridge:} Delivering Louw-Nida \cite{louwnida1988} Semantic Domain Analysis directly into personal prayer and action. \\ \midrule
\textbf{The Hermeneutical Circle Deficit:} AI delivers linear answers, ignoring the iterative part-whole dynamic \cite{gadamer1960}. & Single-pass response with no recursive refinement. & \textbf{Recursive Socratic Cycle:} AI enforces multi-turn dialogue that mimics the hermeneutical spiral \cite{osborne2006}. \\ \bottomrule
\end{tabular}
\end{table*}

\subsection{Addressing the ``Mirror Effect''}
Existing GenAI implementations often suffer from a failure to ``stretch'' the user's cognitive horizon. If a user observes that ``Paul was kind,'' a standard LLM often responds with ``Yes, Paul's kindness is evident.'' The Scriptorium Framework enforces a ``Non-Mirroring'' policy, where the AI must identify the linguistic root of the specific kindness mentioned (e.g., \textit{charis} vs. \textit{eleos} vs. \textit{chrestotes}) and provide an external historical stake to challenge the user's initial shallow observation.

\subsection{The Hermeneutical Spiral}
Osborne \cite{osborne2006} proposes the ``hermeneutical spiral'' as a refinement of the traditional hermeneutical circle: rather than moving in a closed loop, the interpreter spirals upward toward greater understanding with each pass through the text. The Scriptorium Cycle (Section 4) operationalizes this spiral by ensuring that each AI interaction adds a new layer of depth---linguistic, historical, or relational---that was absent in the previous turn. This recursive deepening is what distinguishes the framework from a simple question-and-answer chatbot.

\subsection{Lectio Divina as a Precursor}
The Scriptorium Framework also draws inspiration from the ancient practice of \textit{Lectio Divina}, formalized by the 12th-century Carthusian monk Guigo II \cite{guigo1150}. This four-stage method---\textit{Lectio} (Reading), \textit{Meditatio} (Meditation), \textit{Oratio} (Prayer), and \textit{Contemplatio} (Contemplation)---anticipates the Scriptorium's emphasis on moving from analytical observation to relational resting. Where \textit{Lectio Divina} relies on the Holy Spirit as the interpretive guide, the Scriptorium Framework supplements this with computational linguistic analysis, ensuring that the user's contemplation is grounded in textual rigor rather than subjective intuition alone.
